\documentclass{article}%
\usepackage[T1]{fontenc}%
\usepackage[utf8]{inputenc}%
\usepackage{lmodern}%
\usepackage{textcomp}%
\usepackage{lastpage}%
\usepackage{geometry}%
\geometry{margin=2.2cm}%
%
\usepackage{amsmath}%
\usepackage{amssymb}%
\usepackage{booktabs}%
\usepackage{longtable}%
\usepackage{array}%
\usepackage{hyperref}%
\hypersetup{hidelinks}%
\title{IME 2026 {-} Química (claude\_opus\_5\_5) x Poliedro\\[4pt]\large Comparação com o gabarito Poliedro Resolve\thanks{Gerado com gpt-resolve: \url{https://github.com/lgabs/gpt-resolve}}}%
\date{10/11/2026}%
%
\begin{document}%
\normalsize%
\maketitle%
\section*{Resumo}\textbf{Nota média: 10,0/10} (18 de 18 itens com a resolução correta; cada item vale o mesmo dentro da questão).\par\medskip%
\begin{longtable}{@{}c c p{0.78\textwidth}@{}}\toprule Questão & Nota & Resumo \\ \midrule
1 & 10,0 & A resolução está correta e bate com o gabarito do Poliedro. Os dois chegam a Q = 9648 C (0,1 mol de e$^{-}$), 0,05 mol de M pela relação M$^2$? + 2e$^{-}$ $\to$ M, e massa molar de 87,6 g/mol, que identifica o metal como estrôncio (Sr). Refiz a conta e conferi o enunciado no PDF original, e não há divergência. A resolução ainda exclui Rb e Y e confirma que o Sr forma SrCl$_{2}$. Os dois pontos de dúvida são de baixa importância e não afetam a nota. \\
2 & 10,0 & O item a coincide: C é um álcool terciário. No item b, B (2{-}metilbut{-}1{-}eno) e C (2{-}metilbutan{-}2{-}ol) também coincidem. A única divergência é A: a resolução dá o 1{-}cloro{-}2{-}metilbutano e o gabarito Poliedro dá o 2{-}cloro{-}2{-}metilbutano. Os dois formam B por desidro{-}halogenação e o enunciado não fixa a base nem exige o produto principal, então as duas respostas são válidas. A da resolução é mais rigorosa, porque dá B como único alceno; com o cloreto terciário e KOH/EtOH, B é o produto minoritário (Saytzeff). Não encontrei erro de química na resolução. Os pontos para o professor são a ambiguidade de A e a hipótese, comum às duas soluções, de que não há rearranjo de carbocátion na hidratação. \\
3 & 10,0 & A resolução está certa nos dois itens e bate com o gabarito do Poliedro. a) O acetato de linalila é C12H20O2: H2C=CH--C(CH3)(OCOCH3)--CH2--CH2--CH=C(CH3)2, ou acetato de 3,7{-}dimetilocta{-}1,6{-}dien{-}3{-}ila. b) O polímero vem da poliadição na dupla trissubstituída C6=C7, com unidade de repetição --{[}CH(R)--C(CH3)2{]}n-- e R = --CH2CH2C(CH3)(OAc)CH=CH2. É a mesma unidade do gabarito, só escrita na ordem inversa. Refiz as contas e as fórmulas conferem: 20 H, 3 insaturações, M $\approx$ 196,3 g/mol, e a unidade de repetição tem a fórmula do monômero. A transcrição do enunciado confere com a prova original. Não encontrei erro nem divergência relevante. \\
4 & 10,0 & A resolução está correta e bate com o gabarito do Poliedro: 7,2 kg (7200 g) de ácido acético por piscina. O valor vem de C = 6,0$\times$10$^{-}$$^{4}$ mol/L e n = 120 mol. A resolução usa balanço de cargas e de massa, e o gabarito usa a Lei de Ostwald sem aproximação ($\alpha$ = 1/6); os dois métodos são equivalentes. A única diferença é 7206 g contra 7200 g, que vem de usar 60,05 ou 60 g/mol na massa molar. \\
5 & 10,0 & A resolução e o gabarito do Poliedro concordam: m1/m2 = 2$\cdot$M(CH4)/M(O2) $\approx$ 1. A resolução dá 1,003 porque usa as massas molares completas da tabela, e o gabarito dá exatamente 1 porque usa 16 e 32. Isso é só arredondamento. O raciocínio está certo nos dois: o $\Delta$Hc se cancela, a primeira queima tem avanço 4 vezes maior que a segunda e cada mol de CH4 consome 2 mol de O2. Conferi o enunciado no PDF original e as contas em python, e os dois estão certos. \\
6 & 10,0 & A resolução está correta e bate com o gabarito do Poliedro: os dois chegam ao ácido etanoico (ácido acético), C2H4O2 (CH3--COOH, ou H3C--C(=O)--OH). Refiz as contas e elas fecham: M = 30$\cdot$2 $\approx$ 60 g/mol, 3,0$\times$10$^{-}$$^3$ mol de NaOH para 3,0$\times$10$^{-}$$^3$ mol de ácido (um único H ionizável), e C2H4O2 é a única fórmula possível. Os caminhos são equivalentes: o gabarito resolve 12x + y = 28 e a resolução usa M(R) = 15, ou seja, CH3. Ficaram só observações menores sobre hipóteses do enunciado (dissociação completa e vapor sem dímeros) e o arredondamento de 60,5 para 60, que não mudam a nota. \\
7 & 10,0 & A resolução está correta e coincide com o gabarito Poliedro: mínimo de 42 placas. Os valores intermediários também batem: T\_fus = 1296,48 K, $\Delta$H\_fus $\approx$ 12,96 kJ/mol, $\Delta$H\_global $\approx$ +49,93 kJ por mol de malaquita, E $\approx$ 83,8 MJ e 41,91 placas, arredondado para 42. Refiz as contas em python e todas conferem. O único ponto a discutir com o professor é a unidade 'kJ/mol' das expressões de G no enunciado, que não é compatível com $\Delta$S\_fus = 10 J/(mol$\cdot$K). Ambos contornam isso da mesma forma e o resultado não muda. \\
8 & 10,0 & A resolução e o gabarito do Poliedro Resolve concordam e os dois estão certos: a deficiência é de cerca de 40,5\%. O ânion tem 50 elétrons de valência. Se todas as ligações fossem comuns de 2 centros e 2 elétrons, seriam 30 ligações B--B (as arestas do icosaedro) mais 12 B--H, ou seja, 42 ligações e 84 elétrons. Faltam 34, e (84 - 50)/84 $\approx$ 40,5\%. Conferi as contas com python3 e o enunciado no PDF original. A resolução ainda traz duas verificações corretas que o gabarito não tem: o mesmo critério aplicado ao diborano dá 25\%, e a regra de Wade explica os 26 elétrons da gaiola. Não há erro em nenhum dos dois. \\
9 & 10,0 & A resolução está certa e coincide com o gabarito do Poliedro nos três itens: a) {[}CO2{]} $\approx$ 1,8$\times$10$^{-}$$^{5}$ mol/L; b) pH $\approx$ 5,56, que o gabarito arredonda para 5,6 a partir de 2pH = 11,11; c) o pH diminui, porque a ionização é endotérmica ($\Delta$H$^\circ$ = +7,6 kJ/mol). Refiz as contas em Python e todas conferem ($\Delta$G$^\circ$ = +36,3 kJ/mol, ln K = -14,68, pH = 5,557, $\Delta$H$^\circ$ = +7,6 kJ/mol). As únicas ressalvas são menores: usar k\_H de 300 K para 25 $^\circ$C e supor P = 1 atm. Nenhuma diferença afeta a nota. \\
10 & 10,0 & A resolução da Q10 bate com o gabarito Poliedro nos três itens, e os dois estão corretos. Em (a), a sequência é $\beta$$^{-}$, $\alpha$, $\alpha$, $\alpha$ (Ac $\to$ Th $\to$ Ra $\to$ Rn $\to$ Po), e a resolução ainda dá a equação global balanceada. Em (b), a série é a do actínio (4n+3), com uma demonstração melhor que a do gabarito. Em (c), $\lambda$ $\approx$ 0,173 s$^{-}$$^{1}$ e $\tau$ $\approx$ 5,77 s, valores que conferi em python. A nota do gabarito sobre '277' no enunciado não se confirma no PDF original, que traz 227. \\
\bottomrule\end{longtable}%
\section*{Principais pontos de dúvida}%
\begin{itemize}%
\item%
\textbf{Q2} \textit{(statement)} O composto A não fica determinado de forma única. Tanto o 1{-}cloro{-}2{-}metilbutano (resolução) quanto o 2{-}cloro{-}2{-}metilbutano (gabarito) formam 2{-}metilbut{-}1{-}eno por desidro{-}halogenação. Com o segundo, o 2{-}metilbut{-}1{-}eno só é o produto principal se a base for volumosa (Hofmann); com KOH/EtOH ele é minoritário (Saytzeff). Vale confirmar com o professor ou com o critério oficial se as duas estruturas de A são aceitas. Se o corretor exigir a estrutura do gabarito, a nota do item b muda.%
\end{itemize}%
\section*{Questão 1}%
\href{https://poliedroresolve.sistemapoliedro.com.br/vestibulares/ime/2026/ime-2-fase-dia-3-25-10-2025-2-fase-dia-3/dissertativa-1-quimica-2-fase-dia-3-ime-2026}{Resolução do gabarito}\par\smallskip%
\begin{longtable}{@{}p{0.06\textwidth} p{0.33\textwidth} p{0.33\textwidth} p{0.1\textwidth} p{0.08\textwidth}@{}}\toprule Item & Resolução & Gabarito & Veredito & Certo \\ \midrule
Q1 -- identificação do metal M & Q = 9648 C = 0,100 mol e$^{-}$; M$^2$? + 2e$^{-}$ $\to$ M; n(M) = 0,050 mol; massa molar = 87,6 g/mol $\Rightarrow$ estrôncio (Sr, Z = 38), sal SrCl$_{2}$ & Q = 0,4$\cdot$6,7$\cdot$3600 = 9648 C = 0,1 F; M$^2$? + 2e$^{-}$ $\to$ M; massa molar = 87,6 g/mol $\Rightarrow$ estrôncio (Sr) & igual & ambos \\
\bottomrule\end{longtable}%
\begin{itemize}%
\item%
\textbf{Q1 -- identificação do metal M)} Refiz a conta: 0,4 $\times$ 6,7 $\times$ 3600 = 9648 C; 9648/96480 = 0,1 mol e$^{-}$; 0,1/2 = 0,05 mol de M; 4,38/0,05 = 87,6 g/mol, que bate com o Sr (87,62) da tabela fornecida. O enunciado da transcrição confere com o PDF original (página 2). A resolução e o gabarito seguem o mesmo raciocínio e chegam à mesma conclusão. A resolução também verifica Rb (85,47) e Y (88,91) e confirma que o Sr forma SrCl$_{2}$.%
\end{itemize}%
\paragraph{Pontos de dúvida}%
\begin{itemize}%
\item%
\textit{(solution)} A resolução supõe eficiência de corrente de 100\% e corrente constante. O enunciado não diz isso, mas é a hipótese padrão, e a resolução a declara.%
\item%
\textit{(solution)} Na verificação, a resolução diz que o Sr metálico é obtido 'justamente por eletrólise do sal fundido'. A eletrólise ígnea de SrCl$_{2}$ (com KCl) existe, mas boa parte da produção industrial usa redução aluminotérmica do SrO. É só um comentário qualitativo e não muda a resposta.%
\end{itemize}%
\section*{Questão 2}%
\href{https://poliedroresolve.sistemapoliedro.com.br/vestibulares/ime/2026/ime-2-fase-dia-3-25-10-2025-2-fase-dia-3/dissertativa-2-quimica-2-fase-dia-3-ime-2026}{Resolução do gabarito}\par\smallskip%
\begin{longtable}{@{}p{0.06\textwidth} p{0.33\textwidth} p{0.33\textwidth} p{0.1\textwidth} p{0.08\textwidth}@{}}\toprule Item & Resolução & Gabarito & Veredito & Certo \\ \midrule
a & Álcool (álcool terciário, que não é oxidado pelo KMnO4/H+) & Álcool (álcool terciário, pois não reage com KMnO4 em meio ácido) & igual & ambos \\
b -- estrutura de A & 1{-}cloro{-}2{-}metilbutano, ClCH2--CH(CH3)--CH2--CH3 & 2{-}cloro{-}2{-}metilbutano, CH3CH2C(CH3)2Cl & diferente & ambos \\
b -- estrutura de B & 2{-}metilbut{-}1{-}eno, H2C=C(CH3)--CH2--CH3 & 2{-}metilbut{-}1{-}eno, CH3CH2C(CH3)=CH2 & igual & ambos \\
b -- estrutura de C & 2{-}metilbutan{-}2{-}ol, (CH3)2C(OH)--CH2--CH3 & 2{-}metilbutan{-}2{-}ol, CH3CH2C(CH3)2OH & igual & ambos \\
\bottomrule\end{longtable}%
\begin{itemize}%
\item%
\textbf{a)} As duas respostas são iguais. O KMnO4 em meio ácido oxida álcoois primários e secundários porque o carbono da hidroxila tem H; o álcool terciário não tem esse H e por isso não reage. A função pedida é álcool.%
\item%
\textbf{b -- estrutura de A)} Refiz o raciocínio. B é o 2{-}metilbut{-}1{-}eno e há dois cloretos C5H11Cl que o formam por desidro{-}halogenação. (i) 1{-}cloro{-}2{-}metilbutano: o único carbono $\beta$ é o C2, que tem um só H, então a eliminação forma só a dupla C1=C2 e B é o único alceno. (ii) 2{-}cloro{-}2{-}metilbutano: os carbonos $\beta$ são C1 (CH3) e C3 (CH2). Com KOH/EtOH, pela regra de Saytzeff, o produto principal é o 2{-}metilbut{-}2{-}eno e B sai como produto minoritário. Com base volumosa (t{-}BuOK, produto de Hofmann), B passa a ser o principal. Nos dois casos A sofre desidro{-}halogenação e forma B, como diz o enunciado. As duas fórmulas têm 11 H, o que confere com C5H11Cl. O enunciado não especifica a base nem pede o produto principal, então as duas respostas são quimicamente válidas. A da resolução é mais rigorosa, porque dá B sem outro alceno. A do gabarito é a escolha didática mais comum, já que o Cl fica no carbono que recebe o OH em C. A banca provavelmente aceita as duas.%
\item%
\textbf{b -- estrutura de B)} É a mesma estrutura. A ozonólise dá HCHO e butanona, e há formaldeído porque a dupla é terminal. Entre os três alcenos C5H10 com dupla terminal, só este dá álcool terciário por hidratação Markovnikov direta. A resolução e o gabarito fazem a mesma triagem.%
\item%
\textbf{b -- estrutura de C)} É a mesma estrutura: um álcool terciário, C5H12O, produto Markovnikov da hidratação de B.%
\end{itemize}%
\paragraph{Pontos de dúvida}%
\begin{itemize}%
\item%
\textit{(statement)} O composto A não fica determinado de forma única. Tanto o 1{-}cloro{-}2{-}metilbutano (resolução) quanto o 2{-}cloro{-}2{-}metilbutano (gabarito) formam 2{-}metilbut{-}1{-}eno por desidro{-}halogenação. Com o segundo, o 2{-}metilbut{-}1{-}eno só é o produto principal se a base for volumosa (Hofmann); com KOH/EtOH ele é minoritário (Saytzeff). Vale confirmar com o professor ou com o critério oficial se as duas estruturas de A são aceitas. Se o corretor exigir a estrutura do gabarito, a nota do item b muda.%
\item%
\textit{(statement)} A resolução e o gabarito descartam o 3{-}metilbut{-}1{-}eno porque usam a adição de Markovnikov sem rearranjo. Na hidratação ácida real, o carbocátion secundário em C2 sofre migração 1,2 de hidreto e vira o carbocátion terciário em C3, o que também leva ao 2{-}metilbutan{-}2{-}ol (álcool terciário). Nesse caso B = 3{-}metilbut{-}1{-}eno também cumpriria as condições do enunciado, e A seria, por exemplo, o 1{-}cloro{-}3{-}metilbutano. No nível do vestibular essa possibilidade costuma ser ignorada. A resolução declara essa hipótese de forma explícita; o gabarito não.%
\item%
\textit{(key)} Na resolução do gabarito, pela descrição da imagem, o 3{-}metilbutan{-}2{-}ol aparece desenhado com uma dupla junto ao carbono do OH. Parece um erro de desenho, e as respostas finais não mudam.%
\item%
\textit{(key)} O gabarito escolhe o 2{-}cloro{-}2{-}metilbutano sem comentar que, com base comum, a regra de Saytzeff favorece o 2{-}metilbut{-}2{-}eno, cuja ozonólise dá acetona e acetaldeído, sem formaldeído. A justificativa da resolução para preferir o 1{-}cloro{-}2{-}metilbutano é mais completa.%
\end{itemize}%
\section*{Questão 3}%
\href{https://poliedroresolve.sistemapoliedro.com.br/vestibulares/ime/2026/ime-2-fase-dia-3-25-10-2025-2-fase-dia-3/dissertativa-3-quimica-2-fase-dia-3-ime-2026}{Resolução do gabarito}\par\smallskip%
\begin{longtable}{@{}p{0.06\textwidth} p{0.33\textwidth} p{0.33\textwidth} p{0.1\textwidth} p{0.08\textwidth}@{}}\toprule Item & Resolução & Gabarito & Veredito & Certo \\ \midrule
a & C12H20O2 (fórmula mínima C6H10O). Estrutura: H2C=CH--C(CH3)(O--CO--CH3)--CH2--CH2--CH=C(CH3)2, que é o acetato de 3,7{-}dimetilocta{-}1,6{-}dien{-}3{-}ila. & C12H20O2 (fórmula mínima C6H10O). Estrutura: CH2=CH--C(CH3)(OCOCH3)--CH2--CH2--CH=C(CH3)2, que é o acetato de 3,7{-}dimetilocta{-}1,6{-}dien{-}3{-}ila. & igual & ambos \\
b & Poliadição na dupla trissubstituída C6=C7: --{[}CH(R)--C(CH3)2{]}n--, com R = --CH2--CH2--C(CH3)(OCOCH3)--CH=CH2. O vinil e o acetato ficam como grupos laterais. & Poliadição só na dupla trissubstituída C6=C7, com o vinil livre: --{[}C(CH3)2--CH(CH2CH2C(CH3)(OCOCH3)CH=CH2){]}n--. & igual & ambos \\
\bottomrule\end{longtable}%
\begin{itemize}%
\item%
\textbf{a)} Mesma estrutura e mesma fórmula molecular. Conferi: a razão 6:10:1 com 2 O dá C12H20O2. São 3 insaturações (2 C=C e 1 C=O), então não há anel. A cadeia principal tem 8 C, com duplas em 1 e 6. As duas metilas da mesma ponta ficam em C7, formando =C(CH3)2. Os carbonos alílicos não terminais são C3 e C5. O C3 é o menos impedido, porque fica ao lado do vinil e não da dupla trissubstituída, e por isso recebe a metila e o acetato. São 20 H e M $\approx$ 196,3 g/mol. Na figura em tikz aparece CH=C(CH3)--CH3, que é só outra forma de escrever =C(CH3)2.%
\item%
\textbf{b)} É a mesma unidade de repetição, só escrita a partir da outra ponta: CH(R)--C(CH3)2 e C(CH3)2--CH(R). Ela tem a fórmula C12H20O2, igual à do monômero, como deve ser numa poliadição. A resolução justifica ainda a via catiônica pelo carbocátion terciário em C7, o que é coerente e não muda a resposta.%
\end{itemize}%
\paragraph{Pontos de dúvida}%
\begin{itemize}%
\item%
\textit{(statement)} O item v ("carbono alílico, não terminal, de menor impedimento espacial") depende de comparar o impedimento de C3 e C5. As duas resoluções escolhem C3, ao lado do vinil e longe das metilas geminais, e essa é a estrutura real do acetato de linalila. Mesmo assim, o critério de impedimento é qualitativo, e o examinador poderia aceitar justificativas diferentes.%
\item%
\textit{(key)} O gabarito escreve "Dieno com duplas não conjugadas 1 e 6 sem carbono alílico terminal", uma redação confusa: C8 e a metila de C7 são alílicos e terminais. O enunciado só pede que o carbono do acetato seja alílico não terminal. Isso não afeta a resposta.%
\item%
\textit{(solution)} A resolução diz que a via catiônica é a mais plausível e que alquenos trissubstituídos polimerizam mal. São comentários a mais, corretos mas fora do que o enunciado pede. Ela também supõe, sem dizer, um encadeamento cabeça{-}cauda regular. O gabarito faz a mesma suposição.%
\end{itemize}%
\section*{Questão 4}%
\href{https://poliedroresolve.sistemapoliedro.com.br/vestibulares/ime/2026/ime-2-fase-dia-3-25-10-2025-2-fase-dia-3/dissertativa-4-quimica-2-fase-dia-3-ime-2026}{Resolução do gabarito}\par\smallskip%
\begin{longtable}{@{}p{0.06\textwidth} p{0.33\textwidth} p{0.33\textwidth} p{0.1\textwidth} p{0.08\textwidth}@{}}\toprule Item & Resolução & Gabarito & Veredito & Certo \\ \midrule
Questão 4 -- massa mínima de ácido acético por piscina & $\approx$ 7,2 kg ($\approx$ 7206 g; C = 6,0$\times$10$^{-}$$^{4}$ mol/L, n = 120 mol, M = 60,05 g/mol) & 7200 g = 7,2 kg ($\alpha$ = 1/6, M = 6$\cdot$10$^{-}$$^{4}$ mol/L, n = 120 mol, massa molar 60 g/mol) & igual & ambos \\
\bottomrule\end{longtable}%
\begin{itemize}%
\item%
\textbf{Questão 4 -- massa mínima de ácido acético por piscina)} O enunciado confere com o original (p. 2). Refiz a conta: com pH = 4, {[}H?{]} = 1,0$\times$10$^{-}$$^{4}$ mol/L e {[}Ac$^{-}${]} $\approx$ 1,0$\times$10$^{-}$$^{4}$ mol/L (a correção por Kw/{[}H?{]} = 10$^{-}$$^{1}$$^{0}$ é desprezível). Então {[}HAc{]} = (10$^{-}$$^{4}$)$^2$/(2$\times$10$^{-}$$^{5}$) = 5,0$\times$10$^{-}$$^{4}$ mol/L e C = 6,0$\times$10$^{-}$$^{4}$ mol/L. Isso dá 120 mol em 2$\times$10$^{5}$ L. Com 60,05 g/mol saem 7206 g; com 60 g/mol saem 7200 g. A resolução faz o balanço de massa e de cargas, e o gabarito usa Ostwald sem aproximação, Ka = M$\alpha$$^2$/(1-$\alpha$), com $\alpha$ = 1/6. As duas contas são equivalentes e chegam a 120 mol e 7,2 kg. A diferença de 6 g vem só da massa molar arredondada.%
\end{itemize}%
\paragraph{Pontos de dúvida}%
\begin{itemize}%
\item%
\textit{(solution)} Se o aluno usar a aproximação {[}H?{]} $\approx$ $\surd$(Ka$\cdot$C), que vale para $\alpha$ pequeno, encontra C = 5$\times$10$^{-}$$^{4}$ mol/L, 100 mol e $\approx$ 6,0 kg. Como $\alpha$ $\approx$ 17\%, essa aproximação não vale. A resolução e o gabarito avisam isso. Vale perguntar ao professor se a banca daria crédito parcial a quem respondesse 6 kg.%
\item%
\textit{(statement)} O enunciado deixa algumas hipóteses implícitas: o sal do metal pesado não altera a acidez (não há hidrólise do cátion), a precipitação não consome nem libera H?, o volume de ácido adicionado é desprezível ($\approx$ 7 L em 200.000 L) e a água começa com pH 7, contribuição desprezível frente a 10$^{-}$$^{4}$ mol/L. A resolução declara essas hipóteses e todas são razoáveis. Elas não mudam o resultado.%
\item%
\textit{(key)} O gabarito começa com Ka = M$\alpha$$^2$, chega a $\alpha$ = 0,2 e só então refaz a conta sem aproximação. Isso é só a forma de apresentar: o resultado final, $\alpha$ = 1/6 e 7,2 kg, está correto. Usar massa molar 60 g/mol em vez de 60,05 g/mol é só arredondamento.%
\end{itemize}%
\section*{Questão 5}%
\href{https://poliedroresolve.sistemapoliedro.com.br/vestibulares/ime/2026/ime-2-fase-dia-3-25-10-2025-2-fase-dia-3/dissertativa-5-quimica-2-fase-dia-3-ime-2026}{Resolução do gabarito}\par\smallskip%
\begin{longtable}{@{}p{0.06\textwidth} p{0.33\textwidth} p{0.33\textwidth} p{0.1\textwidth} p{0.08\textwidth}@{}}\toprule Item & Resolução & Gabarito & Veredito & Certo \\ \midrule
Razão m1/m2 & m1/m2 = 2$\cdot$M(CH4)/M(O2) = 2$\cdot$16,043/31,998 $\approx$ 1,003 $\approx$ 1 (m1 $\approx$ m2) & m1/m2 = 1 (m1 = m2) & equivalente & ambos \\
\bottomrule\end{longtable}%
\begin{itemize}%
\item%
\textbf{Razão m1/m2)} Refiz a conta. Como as duas queimas são a mesma reação, CH4 + 2 O2 $\to$ CO2 + 2 H2O(l), o calor liberado é proporcional ao avanço $\xi$. Então $\xi$1/$\xi$2 = 324/81 = 4. Daí n(CH4)1 = $\xi$1 = 4$\xi$2 e n(O2)2 = 2$\xi$2, o que dá m1/m2 = (4$\xi$2$\cdot$M\_CH4)/(2$\xi$2$\cdot$M\_O2) = 2$\cdot$M\_CH4/M\_O2. Com massas inteiras (16 e 32) isso dá exatamente 1. Com as massas da tabela da prova (16,0426 e 31,998, conferido em python) dá 1,0027. Para conferir com $\Delta$Hc $\approx$ 890 kJ/mol: m1 $\approx$ 5,84 g e m2 $\approx$ 5,82 g. A diferença entre 1,003 e 1 vem só da precisão das massas molares, e a própria resolução conclui que a razão é $\approx$ 1. O gabarito segue o mesmo raciocínio por proporção (o 324 kJ consome 4y mol de O2, que correspondem a 2y mol de CH4, logo x = 2y) e chega ao mesmo resultado.%
\end{itemize}%
\paragraph{Pontos de dúvida}%
\begin{itemize}%
\item%
\textit{(statement)} A expressão "nas mesmas condições" é lida pelas duas resoluções como a mesma reação (combustão completa, água líquida, mesma T e mesma p), com o mesmo $\Delta$Hc por mol. O enunciado não diz isso de forma explícita para a segunda queima, mas é a leitura natural. Sem essa hipótese não daria para resolver, porque o $\Delta$Hc não é fornecido.%
\item%
\textit{(solution)} A resolução dá 1,003 por usar as massas molares completas da tabela (16,043 e 31,998). O gabarito usa 16 e 32 e dá exatamente 1. É só diferença de arredondamento, e a resolução também declara a razão $\approx$ 1. Não deve afetar a nota.%
\end{itemize}%
\section*{Questão 6}%
\href{https://poliedroresolve.sistemapoliedro.com.br/vestibulares/ime/2026/ime-2-fase-dia-3-25-10-2025-2-fase-dia-3/dissertativa-6-quimica-2-fase-dia-3-ime-2026}{Resolução do gabarito}\par\smallskip%
\begin{longtable}{@{}p{0.06\textwidth} p{0.33\textwidth} p{0.33\textwidth} p{0.1\textwidth} p{0.08\textwidth}@{}}\toprule Item & Resolução & Gabarito & Veredito & Certo \\ \midrule
Q6 -- nome e estrutura molecular do ácido & Ácido etanoico (ácido acético), C2H4O2, M $\approx$ 60 g/mol, monoprótico; estrutura H3C--C(=O)--OH & Ácido etanoico (ácido acético), C2H4O2 (M = 60 g/mol, monoácido); estrutura CH3--COOH = H3C--C(=O)--OH & igual & ambos \\
\bottomrule\end{longtable}%
\begin{itemize}%
\item%
\textbf{Q6 -- nome e estrutura molecular do ácido)} O enunciado transcrito confere com o PDF (p. 3). Refiz as contas: M = 30$\cdot$M(H2) = 30$\cdot$2,0158 $\approx$ 60,5 g/mol ($\approx$ 60). n(NaOH) = 0,1$\cdot$0,030 = 3,0$\times$10$^{-}$$^3$ mol. n(ácido) = 0,18/60 = 3,0$\times$10$^{-}$$^3$ mol, então k = 1 e o ácido é monocarboxílico. A equação 12x + y = 28 só tem como solução química x = 2, y = 4 (x = 1 daria y = 16, impossível). Fica C2H4O2, e o único ácido carboxílico com essa fórmula é o CH3COOH. A resolução chega ao mesmo resultado por outro caminho (R = 60 - 45 = 15, ou seja, CH3; ou 14n + 32 = 60, logo n = 2) e confere a resposta com a massa tabelada de 60,05 g/mol. Nome, fórmula e estrutura são iguais aos do gabarito.%
\end{itemize}%
\paragraph{Pontos de dúvida}%
\begin{itemize}%
\item%
\textit{(statement)} A frase "Considerando a dissociação completa" é uma idealização. O ácido acético é fraco (Ka $\approx$ 1,8$\times$10$^{-}$$^{5}$). Na titulação com NaOH, porém, ele reage por completo de qualquer jeito. A hipótese só serve para garantir que cada H da carboxila consome um OH$^{-}$, e as duas resoluções a usam nesse sentido.%
\item%
\textit{(statement)} Na prática, o vapor de ácido acético forma dímeros, o que daria uma densidade aparente maior que 30 vezes a do H2. O enunciado fixa M $\approx$ 60, ou seja, trata o vapor como monômero. A resolução comenta isso; o gabarito, não. Não muda a resposta.%
\item%
\textit{(solution)} Com a massa do H da tabela, M = 30$\cdot$2,0158 $\approx$ 60,47 g/mol. A resolução arredonda para 60 e justifica pela precisão dos dados. O gabarito usa 60 direto. Isso não afeta k = 1 (com M = 60,47 dá k $\approx$ 1,008).%
\end{itemize}%
\section*{Questão 7}%
\href{https://poliedroresolve.sistemapoliedro.com.br/vestibulares/ime/2026/ime-2-fase-dia-3-25-10-2025-2-fase-dia-3/dissertativa-7-quimica-2-fase-dia-3-ime-2026}{Resolução do gabarito}\par\smallskip%
\begin{longtable}{@{}p{0.06\textwidth} p{0.33\textwidth} p{0.33\textwidth} p{0.1\textwidth} p{0.08\textwidth}@{}}\toprule Item & Resolução & Gabarito & Veredito & Certo \\ \midrule
7 --- número mínimo de placas solares em 1 h & 42 placas ($\Delta$H\_global $\approx$ +49,93 kJ por mol de malaquita, ou seja, por 2 mol de Cu; T\_fus = 1296,48 K; $\Delta$H\_fus $\approx$ 12,96 kJ/mol; E $\approx$ 83,8 MJ; 83,8/2 $\approx$ 41,9, arredondado para 42) & 42 placas ($\Delta$H\_global = +49,93 kJ por 2 mol de Cu; T\_fus = 1296,48 K; $\Delta$H\_fus = +12,96 kJ/mol; E $\approx$ 83,82 MJ; 83,82/2 = 41,91, arredondado para 42) & igual & ambos \\
\bottomrule\end{longtable}%
\begin{itemize}%
\item%
\textbf{7 --- número mínimo de placas solares em 1 h)} Refiz as contas em python. Igualando G\_L = G\_S, obtém{-}se 12964,8 - 10T = 0, logo T\_fus = 1296,48 K. Daí $\Delta$H\_fus = T\_fus$\cdot$$\Delta$S\_fus = 1296,48 $\times$ 10 J/mol = 12,9648 kJ/mol. Para a reação global, $\Delta$H = {[}2(-394) + (-242) + 2(12,9648){]} - (-1054) = +49,9296 kJ por mol de malaquita. Com 213360/63,546 = 3357,57 mol de Cu, temos 1678,78 mol de malaquita e E = 83821 kJ $\approx$ 83,82 MJ. Então 83,82/2 = 41,91, e o mínimo inteiro é 42. Com M = 63,5 g/mol, sai 41,94 e o resultado continua 42. Conferi o enunciado no PDF original (pág. 3): bate com a transcrição, inclusive o rótulo 'kJ/mol' das expressões de G. A resolução e o gabarito seguem o mesmo caminho (Hess + $\Delta$H\_fus = T\_fus$\cdot$$\Delta$S\_fus) e chegam aos mesmos valores intermediários e à mesma resposta final.%
\end{itemize}%
\paragraph{Pontos de dúvida}%
\begin{itemize}%
\item%
\textit{(statement)} O enunciado diz que G\_S e G\_L estão em kJ/mol. Nessa leitura, d(G\_L - G\_S)/dT = -10 daria $\Delta$S\_fus = 10 kJ/(mol$\cdot$K) e $\Delta$H\_fus $\approx$ 1,3$\times$10$^{4}$ kJ/mol, o que contradiz o $\Delta$S\_fus = 10,0 J/(mol$\cdot$K) da folha de dados e é fisicamente absurdo. Tanto a resolução quanto o gabarito evitam o problema usando o $\Delta$S\_fus da folha ($\Delta$H\_fus = T\_fus$\cdot$$\Delta$S\_fus). O gabarito faz a conversão sem comentar: '1296,48$\cdot$10 $\Rightarrow$ 12,9648 kJ'. A resolução aponta a inconsistência explicitamente. Isso não muda a resposta, mas vale levar ao professor como possível erro de unidade no enunciado.%
\item%
\textit{(statement)} As expressões de G trazem o termo constante '-2,7$\times$10$^2$', que é atípico em polinômios do tipo SGTE (provavelmente era um termo em T$^2$ com erro de digitação). Ele aparece igual em G\_S e G\_L e se cancela, então não afeta T\_fus nem a resposta.%
\item%
\textit{(solution)} Hipótese implícita, comum à resolução e ao gabarito: o processo inteiro foi tratado como se ocorresse com $\Delta$H\_f constantes (hipótese iii), com H$_{2}$O no estado gasoso. Ninguém contabilizou o calor sensível para aquecer reagentes/produtos de 298 K até $\approx$1296 K, nem possíveis perdas. O enunciado pede o 'mínimo' e dá a hipótese (iii), então a leitura adotada é razoável, mas não está totalmente explícita.%
\item%
\textit{(key)} No gabarito há pequenas imprecisões de redação: '$\Delta$H\_f' é usado para a entalpia de fusão (pode ser confundido com entalpia de formação), as placas são chamadas de 'células voltaicas' e a energia aparece como '83,82 MJ/h'. Nada disso altera o resultado.%
\end{itemize}%
\section*{Questão 8}%
\href{https://poliedroresolve.sistemapoliedro.com.br/vestibulares/ime/2026/ime-2-fase-dia-3-25-10-2025-2-fase-dia-3/dissertativa-8-quimica-2-fase-dia-3-ime-2026}{Resolução do gabarito}\par\smallskip%
\begin{longtable}{@{}p{0.06\textwidth} p{0.33\textwidth} p{0.33\textwidth} p{0.1\textwidth} p{0.08\textwidth}@{}}\toprule Item & Resolução & Gabarito & Veredito & Certo \\ \midrule
Questão 8 (deficiência percentual de elétrons de valência do {[}B12H12{]}2-) & $\approx$ 40,5\% (50 elétrons reais contra 84 esperados com 42 ligações 2c--2e: 30 B--B e 12 B--H; faltam 34 elétrons) & $\approx$ 40,5\% (50 elétrons reais; 84 esperados com 30 ligações B--B, pela relação de Euler, e 12 B--H; (84-50)/84 $\times$ 100\%) & igual & ambos \\
\bottomrule\end{longtable}%
\begin{itemize}%
\item%
\textbf{Questão 8 (deficiência percentual de elétrons de valência do {[}B12H12{]}2-))} Refiz a conta. Elétrons reais: 12$\cdot$3 + 12$\cdot$1 + 2 = 50. Arestas do icosaedro: 12$\cdot$5/2 = 30, e Euler confere (12 - 30 + 20 = 2). Ligações 2c--2e: 30 + 12 = 42, que pedem 84 elétrons. Deficiência: 34/84 = 40,48\% $\approx$ 40,5\%. A resolução chega às 30 arestas contando 5 arestas por vértice e confere por Euler; o gabarito usa Euler direto. O resultado é o mesmo. As verificações a mais da resolução estão certas: o diborano daria 4/16 = 25\%, e a regra de Wade para closo com n = 12 dá 13 pares de esqueleto, ou seja, 50 - 24 = 26 elétrons.%
\end{itemize}%
\paragraph{Pontos de dúvida}%
\begin{itemize}%
\item%
\textit{(statement)} O enunciado pede a deficiência 'em relação ao esperado', e por isso o denominador é 84, como fazem a resolução e o gabarito. Quem dividisse pelo valor real (34/50 = 68\%) ou desse só a diferença absoluta (34 elétrons) chegaria a outra resposta. A leitura adotada é a mais natural do texto.%
\item%
\textit{(statement)} Fica implícito que o 'esperado' corresponde a uma ligação 2c--2e por traço da figura (30 arestas B--B + 12 B--H). Um critério diferente, como completar o octeto de cada B, daria outro número. A resolução e o gabarito adotam a mesma hipótese, que é a coerente com o diborano descrito no enunciado.%
\end{itemize}%
\section*{Questão 9}%
\href{https://poliedroresolve.sistemapoliedro.com.br/vestibulares/ime/2026/ime-2-fase-dia-3-25-10-2025-2-fase-dia-3/dissertativa-9-quimica-2-fase-dia-3-ime-2026}{Resolução do gabarito}\par\smallskip%
\begin{longtable}{@{}p{0.06\textwidth} p{0.33\textwidth} p{0.33\textwidth} p{0.1\textwidth} p{0.08\textwidth}@{}}\toprule Item & Resolução & Gabarito & Veredito & Certo \\ \midrule
a & {[}CO2(aq){]} $\approx$ 1,8$\times$10$^{-}$$^{5}$ mol/L (p\_CO2 = 5,45$\times$10$^{-}$$^{4}$ atm com P = 1 atm; k\_H = 3,3$\times$10$^{-}$$^2$ mol/(atm$\cdot$L)) & {[}CO2{]} $\approx$ 1,8$\times$10$^{-}$$^{5}$ mol/L (P\_CO2 = 545$\times$10$^{-}$$^{6}$ atm; S = k\_H$\cdot$P) & igual & ambos \\
b & pH $\approx$ 5,56 ($\Delta$G$^\circ$ = +36,3 kJ/mol; ln Ka $\approx$ -14,68; Ka $\approx$ 4,3$\times$10$^{-}$$^{7}$; {[}H?{]}$^2$ = Ka$\cdot$C) & pH $\approx$ 5,6 ($\Delta$G$^\circ$ = +36,3 kJ/mol; K = e\^{}(-14,68); {[}H?{]}$^2$ = 1,8$\times$10$^{-}$$^{5}$$\cdot$e\^{}(-14,68); 2pH $\approx$ 11,11) & equivalente & ambos \\
c & O pH diminui: $\Delta$H$^\circ$ = +7,6 kJ/mol (endotérmica), então Ka e {[}H?{]} = $\surd$(Ka$\cdot$C) aumentam com T & O pH diminui: $\Delta$H$^\circ$ = +7,6 kJ/mol, ionização endotérmica, então K e {[}H?{]} aumentam com T & igual & ambos \\
\bottomrule\end{longtable}%
\begin{itemize}%
\item%
\textbf{a)} Mesmo método (Dalton + Henry) e mesmo valor: 3,3$\times$10$^{-}$$^2$$\cdot$5,45$\times$10$^{-}$$^{4}$ = 1,7985$\times$10$^{-}$$^{5}$ $\approx$ 1,8$\times$10$^{-}$$^{5}$ mol/L. A resolução acrescenta uma observação correta: corrigindo k\_H de 300 K para 298 K por van't Hoff ($\Delta$H\_dis = -19,8 kJ/mol), dá k\_H $\approx$ 3,48$\times$10$^{-}$$^2$ e {[}CO2{]} $\approx$ 1,9$\times$10$^{-}$$^{5}$. Isso não muda a resposta principal.%
\item%
\textbf{b)} Os dois fazem o mesmo cálculo. Refazendo: $\Delta$G$^\circ$ = -586,8 - (-386,0 - 237,1) = +36,3 kJ/mol; ln K = -36300/(8,3$\cdot$298) = -14,676; log K $\approx$ -6,37; 2pH = 6,37 + 4,745 = 11,115, então pH = 5,557. O gabarito arredonda para 5,6 e a resolução para 5,56; a diferença é só de arredondamento. Pelo log exato, Ka = 4,23$\times$10$^{-}$$^{7}$; com log K = -6,37, sai 4,27$\times$10$^{-}$$^{7}$, que a resolução escreve como 4,3$\times$10$^{-}$$^{7}$. A diferença é irrelevante. As aproximações da resolução também conferem: {[}H?{]} $\approx$ 2,76$\times$10$^{-}$$^{6}$ e {[}OH$^{-}${]} $\approx$ 3,6$\times$10$^{-}$$^{9}$ mol/L. Tratando o sistema como fechado (C - x no denominador), o pH sai 5,59, como ela afirma.%
\item%
\textbf{c)} $\Delta$H$^\circ$ = -692,0 - (-413,8 - 285,8) = +7,6 kJ/mol > 0. Por van't Hoff, K aumenta com T; com C constante, {[}H?{]} aumenta e o pH cai. A estimativa da resolução, $\Delta$pH $\approx$ -0,02 entre 298 K e 308 K, confere (-0,0216).%
\end{itemize}%
\paragraph{Pontos de dúvida}%
\begin{itemize}%
\item%
\textit{(statement)} A folha de dados dá a constante de Henry a 300 K, mas o item a) pede 25 $^\circ$C (298 K). Resolução e gabarito usam o valor de 300 K direto. A correção por van't Hoff ($\Delta$H\_dis = -19,8 kJ/mol, que se tira dos dados) leva a {[}CO2{]} $\approx$ 1,9$\times$10$^{-}$$^{5}$ mol/L, cerca de 5\% a mais, e muda o pH em só \textasciitilde{}0,01. Vale confirmar com o professor se a banca esperava ou aceitaria essa correção.%
\item%
\textit{(statement)} O enunciado não dá a pressão atmosférica. Os dois supõem P = 1 atm para converter 545 ppm (base molar) em pressão parcial. A hipótese é razoável e está declarada na resolução.%
\item%
\textit{(solution)} No item b), a resolução trata o sistema como aberto ({[}CO2{]} = C fixo, reposto pela atmosfera), o que é coerente com 'em equilíbrio com a atmosfera'. Pelo modelo fechado (C - x), o pH sai 5,59. Os dois arredondam para 5,6, então não muda a nota.%
\item%
\textit{(key)} Erros de digitação no gabarito do Poliedro: a equação do item b) aparece truncada ('H2O + CO2 ? H? + HCO'), e o item c) escreve $\Delta$Hf{[}CO3$^{-}$(aq){]} onde deveria ser HCO3$^{-}$. O valor -692 usado é o do HCO3$^{-}$, então a conta está certa.%
\item%
\textit{(solution)} Ka escrito como 4,3$\times$10$^{-}$$^{7}$ (vem de log Ka = -6,37); pelo ln exato, Ka = 4,23$\times$10$^{-}$$^{7}$. É só arredondamento e não afeta o pH (5,557 a 5,559).%
\item%
\textit{(solution)} No item c), os dois desprezam o aumento de Kw com a temperatura e qualquer efeito sobre a 2ª ionização. Isso é coerente com {[}H?{]} ? {[}OH$^{-}${]} e não altera a conclusão.%
\end{itemize}%
\section*{Questão 10}%
\href{https://poliedroresolve.sistemapoliedro.com.br/vestibulares/ime/2026/ime-2-fase-dia-3-25-10-2025-2-fase-dia-3/dissertativa-10-quimica-2-fase-dia-3-ime-2026}{Resolução do gabarito}\par\smallskip%
\begin{longtable}{@{}p{0.06\textwidth} p{0.33\textwidth} p{0.33\textwidth} p{0.1\textwidth} p{0.08\textwidth}@{}}\toprule Item & Resolução & Gabarito & Veredito & Certo \\ \midrule
a & X = Ac{-}227, J = Th{-}227, L = Ra{-}223, Z = Rn{-}219, T = Po{-}215. Etapas: $^2$$^2$$^{7}$$_{8}$$_{9}$Ac $\to$ $^2$$^2$$^{7}$$_{9}$$_{0}$Th + $^{0}$?$_{1}$$\beta$; $^2$$^2$$^{7}$$_{9}$$_{0}$Th $\to$ $^2$$^2$$^3$$_{8}$$_{8}$Ra + $^{4}$$_{2}$$\alpha$; $^2$$^2$$^3$$_{8}$$_{8}$Ra $\to$ $^2$$^{1}$$^{9}$$_{8}$$_{6}$Rn + $^{4}$$_{2}$$\alpha$; $^2$$^{1}$$^{9}$$_{8}$$_{6}$Rn $\to$ $^2$$^{1}$$^{5}$$_{8}$$_{4}$Po + $^{4}$$_{2}$$\alpha$. Equação global: $^2$$^2$$^{7}$$_{8}$$_{9}$X $\to$ $^2$$^{1}$$^{5}$$_{8}$$_{4}$T + 3 $^{4}$$_{2}$$\alpha$ + $^{0}$?$_{1}$$\beta$. & $^2$$^2$$^{7}$$_{8}$$_{9}$X $\to$ $^2$$^2$$^{7}$$_{9}$$_{0}$J + $^{0}$?$_{1}$$\beta$; $^2$$^2$$^{7}$$_{9}$$_{0}$J $\to$ $^2$$^2$$^3$$_{8}$$_{8}$L + $^{4}$$_{2}$$\alpha$; $^2$$^2$$^3$$_{8}$$_{8}$L $\to$ $^2$$^{1}$$^{9}$$_{8}$$_{6}$Z + $^{4}$$_{2}$$\alpha$; $^2$$^{1}$$^{9}$$_{8}$$_{6}$Z $\to$ $^2$$^{1}$$^{5}$$_{8}$$_{4}$T + $^{4}$$_{2}$$\alpha$ & igual & ambos \\
b & O decaimento $\alpha$ reduz A em 4 e o $\beta$ não altera A, então o resto de A na divisão por 4 se mantém. Como 227, 223, 219 e 215 deixam resto 3 (A = 4n+3), trata{-}se da série do actínio (U{-}235 $\to$ Pb{-}207). & Série (família) do Actínio (4n+3): 227, 223, 219 e 215 deixam resto 3 na divisão por 4; X é o actínio. & igual & ambos \\
c & $\lambda$ = ln2/t$\frac12$ = 0,693/4 $\approx$ 0,173 s$^{-}$$^{1}$; $\tau$ = 1/$\lambda$ = t$\frac12$/ln2 $\approx$ 5,77 s & $\lambda$ = (0,301/0,434)/4 $\approx$ 0,173 s$^{-}$$^{1}$; vida média = 1/$\lambda$ $\approx$ 5,77 s & igual & ambos \\
\bottomrule\end{longtable}%
\begin{itemize}%
\item%
\textbf{a)} As quatro etapas são idênticas: $\beta$$^{-}$ e depois três $\alpha$. A resolução também apresenta a soma global (227 = 215 + 12; 89 = 84 + 6 - 1), que está correta e é compatível com as etapas do gabarito. Pelo PDF original, os números de massa de X e J são 227, não 277, então não há discrepância no enunciado.%
\item%
\textbf{b)} A conclusão é a mesma: série do actínio, 4n+3. A resolução vai além do gabarito ao demonstrar por que o resto na divisão por 4 não muda, o que atende melhor ao verbo 'demonstre'. Conferi: 227, 223, 219, 215, 207 e 235 deixam todos resto 3 na divisão por 4.%
\item%
\textbf{c)} Refiz as contas: 0,301/0,434 = 0,6935, então $\lambda$ = 0,1734 s$^{-}$$^{1}$ e $\tau$ = 4/0,6935 = 5,767 s $\approx$ 5,77 s. Com o valor exato de ln 2, $\tau$ = 5,771 s. As duas respostas coincidem, e as unidades (s$^{-}$$^{1}$ e s) estão no SI.%
\end{itemize}%
\paragraph{Pontos de dúvida}%
\begin{itemize}%
\item%
\textit{(key)} A nota do gabarito diz que o enunciado traz 277 para X e J. O PDF original (p. 5) mostra claramente $^2$$^2$$^{7}$$_{8}$$_{9}$X e $^2$$^2$$^{7}$$_{9}$$_{0}$J, então essa observação não procede para a prova oficial. Ela não muda nenhuma resposta.%
\item%
\textit{(key)} O gabarito escreve 1/0,173 = 5,77 s, mas 1/0,173 = 5,780. O valor 5,77 vem de 4/0,6935 = 5,767. É só arredondamento intermediário, sem efeito na nota.%
\item%
\textit{(statement)} O enunciado usa a letra Z para nomear o nuclídeo, e a resolução também usa Z para o número atômico ($\Delta$Z, 'Z = 89'). A troca é só de notação e não gera erro, mas pode confundir quem corrige.%
\item%
\textit{(statement)} 'Equação estequiométrica da série' admite duas leituras: as etapas separadas ou uma equação global. A resolução apresenta as duas formas e o gabarito só as etapas. Qualquer uma deve ser aceita.%
\end{itemize}%
\end{document}